% !TeX root = ../main.tex

\chapter{Introduction}
\label{ch:introduction}

\section{Weather Forecasting and Its Societal Stakes}
\label{sec:intro-motivation}

Accurate short-range weather forecasting is a public-safety service. In Taiwan, a small, densely populated island lying along the western Pacific typhoon track, the stakes are especially concentrated: a single landfalling typhoon can dump more than \SI{1000}{\milli\metre} of rainfall on steep orography in twelve hours, trigger flash floods and landslides, and disrupt transportation, agriculture, and power. Between typhoons, Mei-yu frontal systems, afternoon convective storms, and orographically forced rainfall drive the remainder of the hydrological cycle. For decision-makers who must issue evacuation orders, schedule reservoir releases, or re-route air traffic, the relevant question is almost always probabilistic and local: \emph{how much rain, where, within the next six to twelve hours?}

The operationally critical quantity --- precipitation at convective scales --- is also the hardest for weather models to predict. Precipitation fields are sparse (most pixels are dry), heavy-tailed (the few wet pixels span orders of magnitude), and strongly coupled to terrain and mesoscale dynamics that only emerge at kilometre-scale resolution. This combination means that any forecasting system which produces visually smooth outputs, or which regresses toward the mean, will systematically underpredict the very events that matter most.

\section{A Brief History of Weather Prediction Methods}
\label{sec:intro-history}

Modern weather forecasting has passed through four broad eras, each driven by a different scientific or computational advance.

\paragraph{Numerical Weather Prediction (NWP).} Since Richardson's early twentieth-century vision and its post-war realisation by Charney, Phillips, and others, operational forecasting has been dominated by numerical integration of the primitive equations of atmospheric motion. State-of-the-art systems such as the ECMWF Integrated Forecasting System (IFS), the NCEP Global Forecast System (GFS), and regional convection-permitting models like WRF and its Taiwan variant RWRF discretise the atmosphere on a three-dimensional grid, solve the governing fluid-dynamics and thermodynamics equations, and close the unresolved sub-grid processes (convection, microphysics, boundary-layer turbulence, radiation) through physical parameterisations. Uncertainty is quantified by ensemble prediction systems that perturb initial conditions and model physics. NWP is physically consistent, well-understood, and remarkably skilful, but convection-permitting runs are computationally expensive, parameterisation error is non-trivial, and operational latency is a fundamental constraint.

\paragraph{Statistical and classical machine-learning post-processing.} Long before deep learning, statistical techniques --- model output statistics (MOS), analog methods, kriging, and more recently random forests and gradient-boosted trees --- have been used to bias-correct and downscale NWP outputs. These approaches treat the NWP forecast as a feature vector and learn a local correction to each observation site. They set an important precedent: that learned components can complement physical cores, and that the ultimate criterion is predictive skill at points of interest rather than global physical fidelity.

\paragraph{The deep-learning era.} Starting in the early 2020s, a new generation of purely data-driven global weather models began to rival and in some cases surpass operational NWP at synoptic scales. FourCastNet \citep{pathak2022fourcastnet} used an Adaptive Fourier Neural Operator to produce six-hourly global forecasts at \ang{0.25} resolution. Pangu-Weather \citep{bi2023pangu} introduced a three-dimensional hierarchical transformer and demonstrated medium-range skill competitive with IFS. GraphCast \citep{lam2023graphcast} adopted an icosahedral-mesh graph neural network and achieved operational-quality ten-day forecasts at a fraction of the inference cost of IFS. At convective scales, MetNet \citep{sonderby2020metnet,espeholt2022metnet2,andrychowicz2023metnet3} and the Deep Generative Model of Rainfall (DGMR) \citep{ravuri2021dgmr} pushed radar-based nowcasting from deterministic CNNs toward generative models capable of producing sharp, ensemble-like outputs. Most recently, StormCast \citep{pathak2024stormcast} demonstrated that a convection-permitting diffusion model can learn a km-scale regional forecaster directly from reanalysis data.

\paragraph{The generative turn.} A persistent problem of deterministic deep-learning forecasters is that, because they are trained with pixelwise regression losses, they produce fields that are systematically blurry and that underestimate extremes. Generative models resolve this by learning the \emph{distribution} of plausible forecasts rather than its conditional mean. DGMR showed this qualitatively for radar nowcasting; GenCast \citep{price2024gencast} made the argument at global scale, demonstrating that a conditional diffusion model produces both sharper and better-calibrated ensemble forecasts than GraphCast. The lesson transfers directly to regional convective forecasting, and it motivates the generative forecaster at the centre of this thesis.

\section{Diffusion Models for Weather: Intuition}
\label{sec:intro-diffusion}

A \emph{diffusion model} is a generative model trained to reverse a gradual noising process. Given a clean sample $x_0$ drawn from some data distribution --- say, a high-resolution precipitation field over Taiwan --- we progressively add Gaussian noise until after enough steps the result is indistinguishable from pure noise. A neural network is then trained to undo one step of this process: given a noised sample and the current noise level, predict the underlying clean signal (equivalently, the score of the noised distribution). At sampling time we start from pure Gaussian noise and apply the learned denoiser repeatedly until we recover a clean sample.

Three properties make this framework attractive for weather forecasting. First, it is naturally \emph{probabilistic}: different noise seeds produce different but equally plausible samples, giving an ensemble for free. Second, the samples are \emph{sharp}: because the model is never asked to average over outcomes, it does not regress toward the mean. Third, the denoiser can be \emph{conditioned} on auxiliary inputs --- the previous atmospheric state, coarse boundary fields, invariant land-surface information --- which is exactly the interface a forecaster needs. GenCast at global scale and StormCast at regional scale both exploit these properties.

\section{The Inference-Cost Problem}
\label{sec:intro-cost}

The cost of this flexibility is paid at inference time. High-quality diffusion sampling requires integrating the model's \emph{probability-flow ordinary differential equation} (PF-ODE) with a few tens of neural-network forward passes per generated sample --- the so-called \emph{number of function evaluations} (NFE); the 18-step Heun schedule used by the baseline in this thesis costs 35 NFE per forecast step. For a single forecast step this is tolerable, but operational use cases compound it. An autoregressive 24-hour rollout multiplies the cost by 24. A 50-member ensemble for probabilistic forecasting multiplies it again. A convection-permitting Taiwan-domain forecast with a 50-member ensemble over a 24-hour horizon therefore requires on the order of $50 \times 24 \times 35 = 42{,}000$ forward passes of a large U-Net --- which places the entire class of generative forecasters near the edge of, or outside, the operational envelope of a regional forecasting centre that must refresh its nowcast every few minutes between radar volumes.

The cost is not the only issue. The reason diffusion needs so many steps is that its PF-ODE trajectory --- the path that carries a pure-noise sample to a clean forecast --- is \emph{curved}. A curved trajectory cannot be integrated accurately in one or two steps; coarse step counts incur visible discretisation error, and that error is worst precisely on the sparse, heavy-tailed precipitation channel, where the score of the noised distribution is ill-conditioned near small noise levels. The practical symptom is that an aggressively sub-sampled diffusion residual reverts toward a smooth, ``slightly wet everywhere'' field that scores acceptably on RMSE but destroys the convective structure operational users care about.

\section{From Diffusion to Flow Matching}
\label{sec:intro-flow}

Both problems --- the step count and the curvature --- point at the same lever: the \emph{shape of the generative trajectory}. If the path from noise to data were a \emph{straight line} rather than a curved diffusion ODE, then a handful of explicit Euler steps would integrate it almost exactly, and a single step would already be a strong approximation.

\emph{Conditional Flow Matching} (CFM) \citep{lipman2023flow,tong2024improving} is a recently-proposed generative framework that learns exactly such a model. Instead of learning a score and integrating a curved PF-ODE, CFM regresses a deterministic \emph{velocity field} that transports a Gaussian prior to the data along a straight-line probability path, and it does so with a simulation-free mean-squared-error objective that needs \emph{no teacher model at all}. \citet{ribeiro2025flowcast} applied this idea to radar precipitation nowcasting under the name \textbf{FlowCast} and reported that a 1-to-10-step flow-matching sampler matched or beat a 100-step diffusion sampler on probabilistic precipitation skill at a fraction of the inference cost. A straight-line flow is integrated accurately in a handful of explicit Euler steps --- but it is still \emph{integrated}, and that residual handful of network evaluations is the last cost a flow-matching forecaster leaves on the table.

This thesis pushes the same lever one decisive notch further with \textbf{MeanFlow} \citep{geng2025meanflow}. Rather than learning the \emph{instantaneous} velocity of the flow and then integrating it, MeanFlow learns the flow's \emph{average} velocity over a whole time interval, so that one network evaluation transports a noise sample across the entire trajectory --- there is no ODE solver and no discretisation error to amortise, and sampling collapses to one or two network evaluations. We adapt MeanFlow to the regional convection-permitting StormCast setting on the Taiwan RWRF domain as the generative residual head, keeping StormCast's two-stage decomposition (a frozen deterministic regression mean plus a generative residual head) intact and replacing only the residual head's curved EDM diffusion process. MeanFlow is the method of this thesis; FlowCast's straight-line I-CFM velocity field is retained as the natural \emph{flow-matching baseline} that MeanFlow generalises --- in fact MeanFlow's objective reduces \emph{exactly} to FlowCast's whenever the time interval has zero length, so the two differ only by the average-velocity bootstrap and a comparison between them is a controlled A/B on the training objective. The EDM diffusion residual is retained throughout as the \emph{reference baseline} that we set out to accelerate, against which all heads are measured at a matched training-sample budget, so that every quality and cost difference can be attributed to the choice of generative trajectory rather than to data, architecture, or the regression mean, all of which are held fixed. Crucially, MeanFlow is \emph{not} a distillation of the diffusion teacher: it is trained from scratch in a single run, which makes it a clean test of whether a better-shaped generative target alone --- without any teacher signal --- already delivers the one-to-two-evaluation speed and the sharpness that operational regional forecasting needs.

\section{Contributions}
\label{sec:intro-contributions}

This thesis makes the following contributions:
\begin{enumerate}
  \item We bring the \emph{average-velocity} \textbf{MeanFlow} objective \citep{geng2025meanflow} into weather forecasting --- to our knowledge for the first time --- and adapt it to the StormCast two-stage regression-plus-residual architecture as the generative residual head: the network learns the displacement of the flow over whole time intervals via the MeanFlow identity (a single-JVP, stop-gradient bootstrapped target), so the residual is sampled in $1$--$2$ network evaluations with no ODE solver, teacher-free and in a single training run, on a four-channel Taiwan regional target.
  \item We retain two reference heads on the same backbone and budget so that the MeanFlow comparison isolates the generative trajectory: the native \emph{EDM diffusion} residual (the 35-NFE model we set out to accelerate), and \textbf{FlowCast}, the straight-line Independent-CFM velocity baseline \citep{lipman2023flow,tong2024improving,ribeiro2025flowcast} that MeanFlow generalises. Because MeanFlow's objective collapses exactly to FlowCast's at zero interval length, MeanFlow-versus-FlowCast is a controlled A/B test on the training objective alone.
  \item We carry out the engineering needed to make all three residual heads share \emph{exactly} the frozen regression mean and per-channel normalisation: same data, same $\Freg$, same $\sigma_{\text{data}}$, same backbone, only the EDM process versus the I-CFM velocity field versus the MeanFlow average-velocity field.
  \item We exploit the four-channel target with a precipitation-aware loss: a channel-weighted L2 objective plus an auxiliary radial power-spectral-density (log-PSD) regulariser on the sparse, heavy-tailed \texttt{qpepre} channel, and we ablate the precipitation channel weight $w_{\text{pr}}$ and the \texttt{log1p}-versus-raw-mm/h precipitation encoding.
  \item We assemble a precipitation-aware evaluation suite --- per-channel RMSE/MAE, probabilistic CRPS, categorical skill scores (CSI, FSS, HSS, false-alarm ratio) at operational thresholds, radial log-PSD, and 1--12-hour rollout errors --- and run it on the full 2022 hourly validation year of the Taiwan RWRF dataset.
  \item We present a head-to-head comparison of MeanFlow against the EDM diffusion residual and the FlowCast velocity baseline at a matched sample budget, including an NFE/wall-clock-versus-quality trade-off and a comparison against a legacy old-StormCast baseline, on a regional convection-permitting forecaster --- to our knowledge the first study of average-velocity flow matching on the Taiwan convective-scale domain.
  \item Appendix~\ref{app:derivations} provides line-by-line derivations of the theoretical machinery MeanFlow rests on --- the score--denoiser (Tweedie) identity, the VE-SDE / PF-ODE, EDM preconditioning, Conditional Flow Matching, and the MeanFlow average-velocity identity --- in a form that should let a reader with a working knowledge of Gaussian algebra and the Fokker--Planck equation reconstruct every identity from scratch.
\end{enumerate}

\section{Thesis Outline}
\label{sec:intro-outline}

Chapter~\ref{ch:background} surveys weather-prediction methods, a condensed account of the diffusion theory underlying the EDM baseline, and conditional flow matching building up to the average-velocity MeanFlow objective at the centre of this thesis, in enough depth to motivate and define every component used later. Chapter~\ref{ch:method} describes the dataset, the frozen regression mean and the EDM diffusion and FlowCast baselines, and --- as its core --- the MeanFlow average-velocity residual pipeline as implemented in the accompanying codebase. Chapter~\ref{ch:experiments} presents the experimental protocol and the head-to-head results across NFE budgets, precipitation thresholds, and rollout horizons. Chapter~\ref{ch:conclusion} summarises findings, limitations, and future work.
